% Display math whose body is only a quantikz circuit, with hyperref and
% the TikZ library quantikz2 loaded without the package: in \[ \], in a
% numbered equation with a \label and an \eqref, with a \tag, and in a
% theorem. Each case follows a paragraph that names it.
\documentclass{article}
\usepackage{amsmath}
\usepackage{amsthm}
\usepackage{tikz}
\usetikzlibrary{quantikz2}
\usepackage{hyperref}
\newtheorem{theorem}{Theorem}
\begin{document}
Bracket quantikz.
\[ \begin{quantikz} \lstick{$|0\rangle$} & \gate{H} & \qw \end{quantikz} \]

Equation tag quantikz.
\begin{equation}
\begin{quantikz} & \gate{X} & \qw \end{quantikz} \tag{P}
\end{equation}

\begin{theorem}
Theorem quantikz.
\begin{equation}
\begin{quantikz} \lstick{$a$} & \ctrl{1} & \qw \\ \lstick{$b$} & \targ{} & \qw \end{quantikz}
\label{eq:circuit}
\end{equation}
\end{theorem}

Ordinary display.
\[ U = H \otimes H \]

Reference to the circuit: \eqref{eq:circuit}.
\end{document}
